\begin{textsample}{POS Dim 2 – se – Score 31.00 – se/se\_em\_30.txt}  \label{ex:f2_pos_022}
5.8 Sistemáticas de Escolha, Acompanhamento, Avaliação e \textbf{Mobilidade} do Itinerário Para compreender o processo de escolha do percurso \textbf{formativo}, é preciso considerar as condições objetivas e subjetivas dos indivíduos, que abrangem desde as opções \textbf{profissionais} existentes, a disponibilidade de vagas no mercado de trabalho e a relação custo/benefício dos tipos de formação disponíveis em cada área até as pressões sociais, aptidões e desejos de cada um.
Os estudos para \textbf{implantação} de \textbf{cursos} \textbf{Técnicos} e FIC consideram entre os seguintes indicadores : estudos socioeconômicos, planejamento de desenvolvimento econômico regional das APL ; a infraestrutura mínima requerida ; os recursos humanos necessários, potencial econômico das regiões ; as \textbf{vocações} locais ; os \textbf{Itinerários} \textbf{Formativos} e diagnósticos de \textbf{rede} e as demandas socioeconômicas locais. - INFRAESTRUTURA - RECURSOS HUMANOS - POTENCIAL ECONÔMICO - \textbf{CATÁLOGO} NACIONAL DE \textbf{CURSOS} \textbf{TÉCNICOS} - \textbf{ITINERÁRIOS} \textbf{FORMATIVOS} - \textbf{GUIA} DE \textbf{CURSOS} FIC E \textbf{TÉCNICO} - \textbf{PROSSEGUIMENTOS} DE ESTUDOS - \textbf{VOCAÇÕES} LOCAIS - DIAGNÓSTICO DA REDE-DEMANDA - VERTICALIZAÇÃO DE ESTUDOS O portfólio de \textbf{cursos} já em \textbf{oferta} na traz um panorama das possibilidades de \textbf{ofertas} de \textbf{Cursos} \textbf{Técnicos} \textbf{Profissionais}, presenciais e EAD, que contemplam, tanto as formas : integrada, concomitante ; subsequente e FIC ( Formação Inicial Continuada ), pontuado o \textbf{perfil} \textbf{profissional} de conclusão de \textbf{curso} dentro de seus eixos \textbf{estruturantes} e as possibilidades de \textbf{certificação}, sejam elas : intermediária, continuada e verticalizada no \textbf{itinerário} \textbf{formativo}.
Portanto, as expectativas, as condições de laboralidade e as oportunidades de \textbf{carreira} devem orientar o desenvolvimento de \textbf{ofertas} cada vez mais diversificadas e mapear os fatores que influenciam as prioridades e escolhas dos estudantes.
Inclui ainda, a elaboração de acordos de cooperação e parcerias para favorecer a \textbf{oferta} de \textbf{qualidade}, o aproveitamento, o compartilhamento da infraestrutura e das experiências de instituições ; a aproximação com o setor produtivo e a construção de um projeto pedagógico que envolva instituições parceiras.
É importante que as pessoas envolvidas nesse contexto participem da elaboração das estratégias e instrumentos de acompanhamento e avaliação dos processos de ensino e aprendizagem, bem como, do acompanhamento.
ACOMPANHAMENTO E AVALIAÇÃO A Avaliação é uma forma de obter e fornecer informações sobre dimensões do processo educacional, de modo a permitir tomadas de decisões que ensejem mudanças necessárias em busca de um contínuo aperfeiçoamento.
Qualquer que seja o instrumento adotado para a avaliação, deverá ensejar um processo contínuo de acompanhamento, análise e recuperação coletiva das competências constituídas e acumuladas ao longo do processo.
Como a avaliação se alicerça nas competências e habilidades definidas a partir dos \textbf{Perfis} \textbf{Profissionais} de Conclusão do \textbf{Curso}, faz -se necessário, primeiramente, realizar uma avaliação diagnóstica para identificar os níveis das habilidades e competências de cada estudante, do repertório obtido ao longo da sua vida, para que seja aproveitada como ponto de partida para novas aprendizagens.
Deve ser contínua, ( permitindo a reflexão processual das aprendizagens ) e \textbf{formativa}, objetivando a investigação da efetividade de todo sistema desenvolvido, com destaque para o design instrucional desenvolvido para o \textbf{Curso}.
A avaliação caracteriza -se também pelo caráter investigativo e pesquisador e visa acompanhar o percurso da aprendizagem realizada pelo estudante, identificando as experiências e saberes.
Sant’Anna ( 1995 ) aborda em seu estudo sobre a importância da avaliação : > “ Tanto educadores quanto educandos reconhecem o significado de valorar os resultados ou expectativas, seja qual for o aspecto da vida em que estejam envolvidos.
Estamos empenhados em detectar quais as melhores razões que justificam a inclusão da avaliação na instituição escolar, e concluímos : A melhoria da instrução está condicionada a uma avaliação eficiente e eficaz da organização ; o desenvolvimento pessoal só se concretiza se houver parâmetros que incentivem e motivem o processo de crescimento ( p.13/14 ) ” A construção do Itinerário de Formação \textbf{Técnica} e \textbf{Profissional} no \textbf{currículo} de Sergipe envolve as seguintes dimensões de acompanhamento e avaliação : 1.
Em relação às competências e habilidades específicas de formação : - Avaliar o que o indivíduo é capaz de fazer numa determinada situação de aprendizagem e no que ele poderá fazer em decorrência do que aprendeu, a partir de uma visão crítica comprometida com a democracia e a aquisição do saber historicamente acumulado. - Avaliação de Egressos.
Faz -se necessária pesquisa para o acompanhamento dos egressos, com o objetivo de avaliar as condições de trabalho e de renda dos \textbf{profissionais}, o seu campo de atuação \textbf{profissional} no mundo do trabalho, a avaliação de que ele faz da Instituição e do seu \textbf{curso} como egresso, como também a contribuição da formação \textbf{profissional} integrada ao ensino médio na sua formação \textbf{superior} e seu significado em suas vidas, avaliando de certa forma, a contribuição da formação \textbf{profissional} para a continuidade dos estudos. - Acompanhamento do Planejamento anual por instituição educacional. - Avaliação de pares ( gestores, professores, parceiros e estudantes ). 2.
Em relação à dimensão científica – busca de evidências sobre a \textbf{qualidade} dos serviços prestados no campo da formação \textbf{profissional} em relação a : conteúdos, processos ou estratégias escolares, como também uma gama de valores, que gera mudanças cumulativas. \textbf{MOBILIDADE} DO \textbf{ITINERÁRIO} É fundamental que sejam incorporados em todos os \textbf{currículos} de EPT nas instituições educacionais, bem como pelos \textbf{profissionais} da educação responsáveis por sua implementação, elemento de correspondência entre os \textbf{Itinerários} \textbf{Formativos} oferecidos, como por exemplo, as habilidades gerais e específicas relacionadas aos eixos \textbf{estruturantes}.
Desta forma, será mais fácil, inclusive, permitir que estudantes \textbf{cursem} a Formação Geral em uma instituição educacional e o Itinerário \textbf{Formativo} em outra, troquem de Aprofundamentos antes da sua conclusão e mudem de instituição educacional ou \textbf{rede} sem prejuízo para a sua formação, considerando ainda as habilidades específicas de cada \textbf{curso}.
Para facilitar a correspondência e, consequentemente, a \textbf{mobilidade} dos estudantes, deve -se \textbf{alinhar} um formato de ementas e \textbf{carga} \textbf{horária} \textbf{média} dos Aprofundamentos, das \textbf{Eletivas} e do Projeto de Vida.
Assim, um estudante será reconhecido como tal em qualquer lugar, porque as habilidades e competências, além da validade do diploma, serão as mesmas em todos os Centros e instituições educacionais de EPT.
Dessa forma, se um estudante começar o \textbf{curso} em uma cidade e, no meio do processo, mudar -se para outra ou se fizer a formação em um lugar e trabalhar em outro, poderá se transferir sem problemas.
Propomos aqui alguns cenários para a concretização dos \textbf{cursos}, com a integração dos diversos segmentos que atuam com o ensino médio, possibilitando aos estudantes, a partir de uma escolha pessoal, vivenciar o \textbf{itinerário} formação \textbf{técnica} e Profissional.
Entre as possibilidades de \textbf{ofertas}, estão : CENÁRIOS - V Itinerário Integrado ao Ensino Médio - EPT Integrado ao EJA - NEM/Conjunto \textbf{Cursos} FICs - NEM Integrado com EPT Parcial - NEM Concomitante com EPT - EPT Integrado ao EMTI CENÁRIO 01 : EDUCAÇÃO \textbf{PROFISSIONAL}, INTEGRADO À EJA ( EDUCAÇÃO DE JOVENS E ADULTOS ) - É o \textbf{curso} \textbf{técnico}, realizado junto com o ensino médio, nesse caso, o PROEJA, para o atendimento dos estudantes \textbf{adolescentes} a partir de 18 anos, garantindo a utilização de mecanismos específicos para esse tipo de alunado que considerem suas potencialidades, necessidades, expectativas em relação à vida, às culturas juvenis e ao mundo do trabalho ( LDB ).
CENÁRIO 02 : EDUCAÇÃO \textbf{PROFISSIONAL}, INTEGRADO EM TEMPO INTEGRAL - É o \textbf{curso} \textbf{técnico}, realizado junto com o ensino médio das instituições educacionais em tempo integral, ou seja, o \textbf{curso} é realizado em dois turnos.
CENÁRIO 03 : EDUCAÇÃO \textbf{PROFISSIONAL}, FORMA CONCOMITANTE ARTICULADO AO ENSINO MÉDIO - É o \textbf{curso} \textbf{técnico}, realizado paralelo com o ensino médio, podendo ser realizado em instituições parceiras.
CENÁRIO 04 : ENSINO MÉDIO COM UM CONJUNTO DE \textbf{CURSOS} FIC ( FORMAÇÃO INICIAL E CONTINUADA ) – O ensino médio, realizado de forma paralela com o conjunto de \textbf{cursos} de formação inicial e continuada de trabalhadores- FIC.
CENÁRIO 05 : EDUCAÇÃO \textbf{PROFISSIONAL}, INTEGRADA AO ENSINO MÉDIO PARCIAL - É o \textbf{curso} \textbf{técnico}, realizado junto com o ensino médio em um só turno.
Em todas as \textbf{ofertas}, os estudantes vivenciarão, além da BNCC, uma parte Flexível/ diversificada.
E a Base Nacional Comum Curricular busca assegurar aos jovens “ uma formação que, em sintonia com seus percursos e histórias, permita -lhes definir seu projeto de vida, tanto no que diz respeito ao estudo e ao trabalho, como também no que concerne às escolhas de estilos de vida saudáveis, sustentáveis e éticos ”.
Há uma etapa de Preparação básica para o trabalho, comum em todo \textbf{currículo} de educação \textbf{profissional} e \textbf{técnica}, em que se deverá trabalhar os fundamentos de Informática, de \textbf{Legislação} trabalhista, de Ética, Relações Interpessoais, Higiene e Segurança do trabalho aplicados ao \textbf{curso} específico, por meio de situações didáticas que promovam a interação do estudante com vivências reais e intervenções na instituição educacional, na comunidade e em empresas.
Acrescente -se o Projeto de Aprendizagem Interdisciplinar - PAI, I, II, III e IV contribuindo para a compreensão e vivência do estudante dos eixos \textbf{estruturantes}. ( Em Anexo ) E, finalmente, a etapa \textbf{destinada} à Formação \textbf{Técnica} e Profissional, que integra conhecimentos voltados para o mundo do trabalho, o domínio de competências e habilidades \textbf{técnicas} construídas com a prática nos contextos \textbf{profissionais} onde será vivenciada a ocupação escolhida. ( Vide ANEXO III e IV \textbf{cursos} ) Cabe salientar que, é preciso delinear um desenho curricular em que toda a organização do \textbf{curso} permita o aproveitamento contínuo de estudos, de forma \textbf{flexível} e convergente, articulando as competências e habilidades com : as \textbf{cargas} \textbf{horárias}, as estratégias metodológicas, os critérios de avaliação e os recursos didáticos, equipamentos e materiais.
Uma estratégia para definição do desenho curricular é aprender com as experiências entre as instituições educacionais que já \textbf{ofertam} Educação Profissional, de modo que possa identificar e analisar de forma estratégica a organização de \textbf{itinerários} realizada por outras instituições educacionais ou \textbf{redes} de ensino.
Apesar de as instituições educacionais estarem inseridas em realidades e contextos diferentes, apresentando suas próprias especificidades, é importante que haja um diálogo para essa troca de experiências, que podem servir para inspirar a criação, recriação e iniciativas que visem o desenvolvimento e aprimoramento de forma estruturada das unidades instituição educacionais.
Segundo BEVILAQUA ( 2009, p.9 ) sobre a construção coletiva dos \textbf{currículos}, “ É fundamental para isso o desenvolvimento de parcerias, a promoção e a articulação com as demandas de trabalho, a fim de se estabelecer mais amplamente o diálogo entre teoria e prática ”.

% matched lemmas: adolescente, alinhar, carga, carreira, catálogo, certificação, currículo, cursar, curso, destinar, eletivo, estruturante, flexível, formativo, guia, horário, implantação, itinerário, legislação, mobilidade, média, oferta, ofertar, perfil, profissional, prosseguimento, qualidade, rede, superior, técnico, vocação
\end{textsample}
